\section{Evaluation evidence in full, R-1 to R-7}\label{sec:F}

\S\ref{sec:11.1} and \S\ref{sec:11.2} carry the spike, the results table and each requirement's verdict. This appendix
carries the evidence behind those verdicts: the spike's own caveats, the per-requirement narrative,
the occupancy stratification, the baseline and learned-contender comparisons, and the per-seed
detail. Every verdict and every regime qualifier also appears in the body; what is only here is the
working.

\subsection{The spike's two caveats}\label{sec:F.1}

\textbf{Two caveats that the spike's own result file states, and that bound every number above.}
The observer used \emph{true occupancy}, so the healthy residual was pure sensor noise with no
model error; the deployed observer carries the additional error of a smoothed estimate, which makes
0 of 20 an optimistic bound. And the worst healthy seed peaked at \textbf{149.1 ppm against
$h$ = 150 ppm}: a 0.9 ppm margin, with the next-highest seed at 91.6 ppm, so it is a tail draw
rather than a design property. R-3 asks the same statistic to hold over seven days and six rooms,
which is far more samples. That horizon has since been run, and the margin did not survive it:
§\ref{sec:11.2} measures \textbf{0.55} confirmed cases per fault-free week with a \textbf{worst week of 6}
against R-3's limit of one, and measured 1.30 per week on the run this margin predicted. The tuning
is therefore \textbf{no longer merely reported: it is measured, and it fails}. ADR-0010 has since
named why the margin was so thin, and it was not the limit: each CUSUM arm was running against a
plausibility guard that had been corrected onto the measurement, which made the statistic
near-zero-mean and put it in exactly the i.i.d. regime the analytic run length assumes. That is why
$k$ and $h$ are left where the spike put them rather than retuned (§\ref{sec:13.4}).


\subsection{What the observer's occupancy dependency costs, stratified}\label{sec:F.2}

That guard is not free, and Section \ref{sec:13.7} states its cost: in A145 (20 seats) or A115 (10 seats), a
room holding two people has a ceiling far above anything a stuck damper produces, so the fault is
invisible until the room fills, and detection delay therefore depends on occupancy. The harness now
prices that dependency at its two ends rather than describing it: with every room at design
occupancy the p95 confirmation delay is \textbf{6 min}, and on a drawn weekday (the same fault, the
same limit, the same seeds) it is \textbf{257 min}, with 5 of 20 injections never confirmed inside
a twelve-hour window (§\ref{sec:11.2}). Two points are not a curve, but the closed-loop arm turns out to
carry the points between them, and they were in the result file rather than missing from it. Its
20 runs are 20 \emph{different} drawn weekdays, so they did not all present the same head count at
the moment of the first ignored command, and the harness records that head count per run precisely
so the correlation can be checked rather than taken. Stratified by it against A124's design
occupancy of 8, the dependency is monotone and steep: \textbf{nothing at all is confirmed below
0.75 of design occupancy}, and above it the median delay falls from the 122 and 257 min of the two
runs at 0.75, through 77 min over the eight at 0.875, to \textbf{31 min} over the five at 1.00.
§\ref{sec:11.2} gives the table and states plainly what kind of evidence it is: a stratification of the
runs that happened, not a sweep of a controlled input.

\subsection{R-1 to R-7, the measurements behind the verdicts}\label{sec:F.3}


\textbf{R-1 is met on the protocol it specifies and not met in closed loop, and both belong here.}
The stepped arm is what R-1's verification column names: the building settled, one command the
frozen damper ignores, 20 seeds, delay measured from that command. At the deployed $h$ = 150 ppm it
confirms 20 of 20 injections with a p95 of \textbf{6 min} against a 15 min budget, and the sweep is
almost flat: p95 moves from 5 min at $h$ = 50 ppm to 8 min at $h$ = 300 ppm, so a sixfold change in
the limit costs three minutes. That makes $h$ a false-alarm parameter rather than a speed
parameter, which is why §\ref{sec:13.4} prices it against R-3 and not against R-1.
\texttt{fault-monitor}'s own R-1 test reports p95 = 6 min on the same protocol over the same seeds
against a deliberately separate physics implementation, and the two agree to the minute: agreement
across two independent plant implementations being the strongest available evidence that 6 min is a
property of the detector rather than of either simulator.

The closed-loop arm injects the same fault on a drawn weekday with the inner loop commanding and
the supervisor live, and it \textbf{fails}: 15 of 20 detected, median 74 min, \textbf{p95 257 min},
five injections never confirmed inside a twelve-hour window. The cause is observability rather than
detector latency. The CUSUM is driven by the measurement against the plausibility ceiling of
Section \ref{sec:7.2}, the concentration the room would reach at design occupancy, while the controller
issues its first ignored command as soon as the room passes its 800 ppm opening threshold, which on
a drawn weekday happens with far fewer people present than that ceiling assumes. Until the room
fills further the measurement stays under the ceiling and the statistic does not accumulate at all.
The clock R-1 specifies starts at the ignored command; the physics does not start until later. Both
numbers are reported because they answer different questions: 6 min is what the detector does when
handed a step response to judge, and 257 min is what the building hands it.

\textbf{That mechanism is not only argued, it is measured, and the measurement was already in the
result file.} The closed-loop arm's 20 runs are 20 different drawn weekdays, so the room held a
different head count on each when the controller issued its first ignored command, and
\texttt{T-24-detection-delay.csv} carries that head count, \texttt{occupants\_at\_ignored\_command},
for every run. A124's design occupancy is 8 (§\ref{sec:6.1}), so the 20 runs stratify by the
actual-to-design occupancy ratio the plausibility ceiling is sensitive to:

\begin{longtable}[]{@{}
  >{\raggedright\arraybackslash}p{(\columnwidth - 10\tabcolsep) * \real{0.17}}
  >{\raggedright\arraybackslash}p{(\columnwidth - 10\tabcolsep) * \real{0.13}}
  >{\raggedright\arraybackslash}p{(\columnwidth - 10\tabcolsep) * \real{0.11}}
  >{\raggedright\arraybackslash}p{(\columnwidth - 10\tabcolsep) * \real{0.13}}
  >{\raggedright\arraybackslash}p{(\columnwidth - 10\tabcolsep) * \real{0.19}}
  >{\raggedright\arraybackslash}p{(\columnwidth - 10\tabcolsep) * \real{0.27}}@{}}
\toprule
\textbf{Occupants at the ignored command} & \textbf{Ratio to design} & \textbf{Runs} &
\textbf{Confirmed} & \textbf{Delay min / median / max} & \textbf{} \\
\midrule
\endhead
5 & 0.625 & 4 & \textbf{0 of 4} & — & never inside 12 h \\
6 & 0.750 & 2 & 2 of 2 & 122 / — / 257 min & two runs, so no median \\
7 & 0.875 & 8 & 8 of 8 & 45 / \textbf{77} / 92 min & \\
8 & 1.000 & 5 & 5 of 5 & 21 / \textbf{31} / 44 min & \\
— & — & 1 & 0 of 1 & — & no ignored command was ever
issued: the room never passed the opening threshold with the damper below it \\
\bottomrule
\end{longtable}

The dependency is monotone, steep, and it has a floor: \textbf{below 0.75 of design occupancy
nothing is confirmed at all}, because the measurement never reaches a ceiling computed at 1.25
times design. That is the mechanism stated as a measurement rather than as an explanation, and it
is the curve §\ref{sec:7.2}'s marker asked for.

\textbf{What kind of evidence this is, stated rather than left to be assumed.} It is a
stratification of the 20 runs that happened, \emph{not} a sweep of a controlled input, and three
bounds follow from that. The strata are unbalanced (4, 2, 8 and 5 runs), so the 0.75 row is two
observations and carries no median worth printing. The range sampled is only 0.625 to 1.000 of
design, because that is what a drawn weekday puts in an eight-desk office; nothing here says what
happens at 0.4 or at 1.5. And every point is one room, A124, at one limit, $h$ = 150 ppm. A proper
sensitivity curve would hold the occupancy ratio as an input and sweep it, which is a different
experiment and has not been run.

One comparison in the table is worth more than the trend, because it separates two causes that are
easily confused. The stepped arm's p95 of 6 min is \emph{also} at an occupancy ratio of 1.0, and
the closed-loop stratum at the same ratio has a median of 31 min. Those two numbers differ by a
factor of five at identical occupancy, so the residue is \textbf{protocol, not occupancy}: the
stepped arm hands the monitor a settled building, the inner loop switched off and one clean step
to judge, while the closed-loop arm asks it to find the same fault in a room that is still filling
and still being commanded. Occupancy explains the 257 min tail; the protocol explains the gap
between 6 and 31 min at the top of the range.

\textbf{The baseline comparison, which is the project's central claim, needs stating carefully.}
T-25 scores the deployed monitor and a fixed CO₂ threshold on the identical per-minute measured
series from the identical runs: 20 seeds with a stuck damper for detection, the same 20 seeds over a
fault-free week for false alarms. The threshold was tuned over 6 limits $\times$ 5 durations
\emph{on the very runs it is then scored on}, which flatters it, and that is deliberate: the
question is whether a model of what the command owed beats the best a fixed threshold can be made
to look, not whether it beats a careless one.

\textbf{Measured, the monitor is ahead on every axis the experiment reports}: detection
\textbf{0.75} against \textbf{0.45}, alarm volume \textbf{0.55} against \textbf{2.35} per
fault-free week, median delay \textbf{74 min} against \textbf{126}, Youden's J \textbf{$+0.55$}
against \textbf{$+0.30$}, and \textbf{0} healthy rooms flagged against 1. The best the grid can be
tuned to is 900 ppm sustained for 5 min.

\textbf{Until this revision the paragraph above said something weaker, and the correction is worth
owning because it runs in the design's favour.} The prose was describing a \emph{pre}-ADR-0010 run,
in which the monitor sat at 1.30 alarms per fault-free week and the threshold was consequently
ahead on Youden's J at $+0.25$ against $+0.20$: a genuine trade, and the report called it one
rather than claiming a win. ADR-0010 halved the monitor's alarm volume with its detection
unchanged, the committed result file has been post-fix since, and the prose had not caught up: the
report was contradicting its own repository and understating its own result. A correction in this
direction deserves more scepticism than one that costs something, so every figure in the preceding
paragraph is read off \texttt{evaluation/results/T-25-baseline-comparison.csv} as committed, and
re-running T-25 reproduces that file byte for byte.

\textbf{The discriminating comparison is at equal alarm volume}, and it is the one the claim rests
on. At 900 ppm sustained for 30 min the threshold raises \textbf{0.40} alarms per fault-free week
(the closest match in the grid to the monitor's 0.55) and detects \textbf{0.15} of the faults
against the monitor's \textbf{0.75}, taking 312 min to do it against the monitor's 74. So the
honest ratio at equal alarm volume is \textbf{five times}, and the \emph{fifteen}-fold figure this
report used to quote belonged to the pre-ADR-0010 operating point and is not true now. The ratio has
moved twice as this number was re-measured (fifteen, then two and a half, now five), which is
itself the reason it is quoted with the operating point attached rather than on its own.

\textbf{T-29 re-matched the same grid on a \emph{held-out} fault-free week}, and there the threshold
does worse again: \textbf{0.05} detection at 0.65 alarms per week, with a 296 min median delay.
The gap between those two numbers is itself a measurement. Scored on the runs it was tuned on the
threshold reaches 0.15 at a matched alarm volume; scored on a week it has never seen, 0.05.
\textbf{Most of the baseline's apparent strength is the tuning-on-test advantage this comparison
deliberately granted it}, and naming that is worth more than either ratio alone. \textbf{Two bounds still travel with the result}: both
detectors were scored against the generator that produced the faults (§\ref{sec:11.5}), and the monitor also
names the fault and carries its evidence while a threshold reports only that a room is above a
number, which is not a detection-rate advantage and is the reason the system exists. The claim
survives in direction at every one of these points; its magnitude depends on which is quoted, so
all of them are, and \textbf{five is the one to defend}.

\textbf{The third detector in that contest is T-29's learned contender, and it is dropped on the
numbers it lost by.} The test plan specified it as optional and kept only if it wins. It is a
one-class novelty detector fitted on a fault-free commissioning window \emph{only} (it never sees
a labelled fault), scoring each minute by its distance from what the room looked like when it was
healthy, over the features the monitor already computes. Thirty operating points, 20 seeds, scored
on the same held-out week. \textbf{No point on the grid beats the monitor}: to reach its 0.75
detection the contender raises \textbf{11.25} alarms per fault-free week against \textbf{0.55} and
takes a \textbf{518} min median delay against 74: about twenty times the false alarms and seven
times the delay. Loosening it until it detects everything costs 475.10 alarms a week, which is not a
detector.

The interesting part is \emph{why}, and three measurements pin it rather than asserting it. First,
\textbf{the contender's misses are exactly the monitor's}: the same five seeds, at every operating
point but the loosest. Both read the same residual, so what limits detection on this plant is the
evidence the residual carries, and \textbf{a learned decision rule cannot add evidence}. Second,
the contender scores each minute on its own while the CUSUM accumulates; its only lever on
in-control run length is a persistence requirement, which is the one part of it that is hand-set
rather than learned. The observer re-anchors the residual towards each reading, so a stuck damper's
excursion decays rather than persisting: the signal shape a cumulative statistic survives and a
per-sample quantile does not. Third, the loss is not an artefact of its training: all twenty
commissioning windows were clean by the monitor's own judgement, none of the 120 per-room fits was
degenerate, and the alarms it is dropped for are counted on the held-out week. This is the result
the project wanted from an optional contender: not that learning cannot work here, but a measured
statement of what it would have to beat and by how far it failed.

\textbf{R-2 is met on the case it covers, and the case it fails on is worth more than the pass.} On
a drawn weekday the faulted room holds 1.000 of its occupied minutes at or below 1000 ppm on all
five seeds, and the figure is unchanged with the lower CUSUM arm off and with the harness's damper
feedback off, so neither the two-sided statistic nor that fidelity control is doing the work. The
reference arm is what makes the result readable: a run with \emph{no fault at all} scores 1.000 and
meets the criterion 5 of 5 in \emph{every} case, including the one R-2 fails. That is the evidence
that the plant is now correctly sized for the speed it runs at (§\ref{sec:13.3}), and it is what pins the
surviving failure to one thing rather than to the building. With a stuck damper \emph{and} every
room simultaneously at its design head count, the supervised arm holds only 0.022 of occupied
minutes and \textbf{0 of 5 seeds} meet the criterion, against 0.009 with the supervisor dead.
\textbf{All five latch and all five carry a plan} (at minutes 481, 481, 481, 494 and 496), so none
of them fails for want of a diagnosis. They fail for want of air, and the limitation each plan
states says which lever ran out. On \emph{one} seed the solver reports that every zone is faulted,
so no healthy damper is left to trim and the gain of \textbf{1.43} is fan speed alone, 0.70 raised
to its 1.00 ceiling. On another every healthy damper is already at its own minimum, so again only
the fan is left, and there it is the healthy rooms' own 5 \% tolerance rather than the plant that
stops it, at 0.735 and a gain of \textbf{1.05}. On the remaining three the trim runs to 0.82, 0.65
and 0.59 with the fan at 0.77, 0.77 and 0.81, reaching gains of \textbf{1.29, 1.62 and 1.77} with
the healthy-flow deviation spent in full at exactly 0.0500. \textbf{So on four of the five it is
the 5 \% healthy tolerance that binds and not the plant's ceiling}, which is the same bound
§\ref{sec:7.4} names, arriving in the regime where it costs something. So the mechanism works where it runs
and is worth $+0.012$ of occupied minutes, against a fault that costs 0.991. And the gap is
not close: the solver's own comfort rule, $(C_\text{measured}-C_\text{out})/(C_\text{target}-
C_\text{out})$, asks for \textbf{2.3 times} the room's flow at the unsupervised mean of 1754 ppm,
against the \textbf{1.05 to 1.77} the plan reaches. \textbf{Compensation redistributes flow and buys a
bounded amount of new flow from the fan, and a building already at design occupancy needs more than
either gives.} That is the plain statement of R-2's failure and of its mechanism.

\textbf{On every one of those five seeds the trim lever has already been damaged before the
supervisor is asked for anything, and that is the larger finding.} Each run carries at least one
\emph{other} room latched with no fault injected in it: one room on two of the seeds, two on
another and four on a third, and on the one whose plan reports every zone faulted, all five.
A latched room is counted among the faulted ones, so it is no longer a healthy damper the solver may
trim: \textbf{R-3's false alarms consume the lever R-2 depends on}. The same effect appears on the
drawn weekday from the other direction, where on two of the five seeds the room is never confirmed
at all and no plan is reached: R-1's closed-loop miss arriving inside R-2's experiment.
\textbf{R-2's design-occupancy failure is therefore substantially a \emph{consequence} of R-3 and
R-1 rather than an independent shortfall}, which also means it was never going to be fixed in the
solver. ADR-0009 settles the part that could have been: the
healthy-room airflow allowance was swept closed-loop from 0.05 to 0.99, one clean build per
setting, and it does not move R-2's criterion by a single seed (0 of 5 at every setting, with
byte-identical per-seed fractions) because with every healthy damper on
its floor the fan lever is not selective and the plant caps it at 1.43 against the 2.00 to 3.30 the
regime asks. The plant's own bound was then measured, with every unlatched room driven to its
minimum position per room rather than by one common trim and the fan at its ceiling: it reaches
only \textbf{2 of the 5 seeds}, and on a third it \textbf{starves A145 to 1033 ppm}. That last
number refutes the premise the sweep was testing (that surplus air costs fan energy rather than
comfort) because a per-room trim takes air from the open healthy rooms as well as giving it to the
nearly shut ones. The constraint was therefore left at 0.05, with the two sides separated in the
code because they are genuinely different quantities, and R-2's verdict is unchanged (ADR-0009).

\textbf{One finding in this experiment is larger than R-2's verdict, and it is reported here rather
than left in the result file.} Until the final run the solver had a single lever: trim the healthy
dampers so the proportional split hands the frozen room a larger share. On a drawn weekday every
damper sits on its 0.15 occupied-hours floor when a plan is solved, so that lever was exhausted
before it was ever used: every plan came back \texttt{trim 1.00, fan 0.70, gain 1.00} and the
supervisory loop measured $+0.000$ of occupied minutes. At design occupancy the trim did reach three
of the five seeds and bought $+0.008$, and the mean CO₂ over those same runs is what makes that
figure unreadable on its own: the supervised arm's mean in the faulted room was \textbf{1788 ppm
against 1754 ppm with the supervisor dead}, so the mechanism this project is named for left the room
\textbf{34 ppm worse} than no supervision at all, for the price of a fan command, a work order and a
mode change. What found it is the shape of the
experiment rather than a reading of the code: T-26 runs an unsupervised arm beside the supervised
one on the same occupancy draw, and nothing else in this project could have given that comparison a
sign. The fix follows from the split law rather than from tuning: $Q_i = Q_\text{fan}a_i/\Sigma a_j$
is linear in $s$, so raising the fan multiplies \emph{every} room's flow including a frozen room's,
and the same regime now measures \textbf{1618.9 ppm} supervised against 1753.6 unsupervised:
\textbf{135 ppm} of help where there was 34 ppm of harm. Three bounds travel with it. The fan lever is \textbf{not
selective}: with no trim in force $A(t)=A_0$ and both ratios collapse to $r = \gamma = s/s_0$, so the
frozen room's gain equals the healthy rooms' over-supply exactly, which means ADR-0008's 43 \% of
headroom is spendable selectively only \emph{through} a trim and the lever is bounded by the 5 \%
healthy-flow tolerance rather than by the plant's ceiling: the ceiling binds only where no healthy
room is left to over-supply. R-2 still fails this case. \textbf{And the charge this lever once
carried against R-3 is withdrawn rather than renumbered.} It read that spending fan speed moved R-3
from 1.25 to 1.30 confirmed cases per fault-free week, because a false alarm now draws a bounded fan
raise that perturbs all six rooms, with the alternative measured rather than argued: keeping the
no-op trim in the published plan dropped R-3 to 0.55 and would have passed it, while costing R-2
0.054 of occupied minutes and failing it by taking five healthy rooms off CO₂ control to hold them
where they already were. All of those R-3 figures were measured against the \emph{pre}-ADR-0010
detector, in which both plausibility guards were being corrected onto the reading and every CUSUM
arm was running on a filter error. And ADR-0010's attribution accounts for all 26 of those
confirmations without the fan: the 13 cascaded ones ran through the \emph{believed} damper opening
that a confirmation pins in \texttt{AhuFlowModel}, not through the fan command, and the re-anchor
that closes that path is now shipped. ADR-0009 adds that the \emph{size} of the lever does not move
R-3 either, the over-supply allowance having been swept from 0.05 to 0.99 with R-3 unchanged at
every setting. One extra confirmation in twenty weeks is inside the burst distribution those weeks
have, so it is not re-asserted as this lever's price; what is honest is that the lever's separate
cost to R-3 has not been re-measured since the fix, and that the rate with the lever in force is
\textbf{0.55} per week (§\ref{sec:11.2}). The R-2 half of the alternative stands exactly as it was measured.

Case B (a damper genuinely shut at opening 0, which ADR-0008 places outside R-2's scope) is
measured for the bound rather than for the requirement: 0.153 of occupied minutes at or below
1000 ppm, \texttt{compensates()} false, achieved gain exactly 1.000, and the real flow split with
A124 shut and the fan at 70 \% is A125A 16.7, A125 30.7, A116A 49.8, A115 107.3, A124 0.0, A145
208.5 l/s. The simulator cannot express this case: \texttt{VentilationZone}'s invariant clamps a
stuck damper at 0 back to the room's 0.15 floor, which is ADR-0008's own decision rather than a bug.
Case B is therefore integrated from the same \texttt{Co2MassBalance} and the same
\texttt{FlowSplit} the simulator uses, at the same 1 s step and on the same occupancy draw, with the
airflow a shut damper actually receives: zero. The bypass is declared in the harness's own output,
not only here.

\textbf{R-3 is not met, and it is the requirement this design does not meet.} Over 20 fault-free
simulated weeks of all six rooms the monitor raises \textbf{0.55 distinct confirmed cases per week}
against a limit of one: \textbf{16 of the 20 weeks are clean} and the worst still raises 6. Before
ADR-0010 the same experiment measured \textbf{1.30 per week with 9 clean weeks}, so the rate has
more than halved and seven more weeks have come clean while the verdict has not moved: the harness scores
the \emph{worst} week, and that is the criterion this requirement is judged by. The rate is driven
by a minority of weeks rather than spread evenly, which is what a tail-draw mechanism looks like
and is what §\ref{sec:11.1}'s 0.9 ppm margin predicted. A case is counted once, by distinct \texttt{caseId}
reaching \texttt{CONFIRMED} rather than by published events, because a latched case republishes its
state on every probe and counting events would charge the monitor a dozen alarms for one
accusation; the threshold it is compared against is counted the same way, by episode.

\textbf{The mechanism is attributed, and the attribution is worth more than the halved rate.} It
was measured rather than argued: the replay was instrumented and all \textbf{26 pre-fix} false
confirmations (the 1.30 per week over 20 weeks that the sentences above quote, before ADR-0010)
classified by room, hour, CUSUM arm and whether any room had already been pinned. The committed
replay, post-fix, carries \textbf{11}. \textbf{13 were
primary} and \textbf{13 cascaded}; \textbf{17 came from the upper arm} and 9 from the lower; all
six rooms confirmed at least once, at hours spread across the whole working day, so it is neither a
morning fill nor a warm-up artefact. The cause is \textbf{one defect with two faces}.
\texttt{ZoneObserver.correctTowards} pulled \emph{both} plausibility guards 30 \% of the way onto
every reading, so neither guard was a bound: each was a first-order filter of the measurement
sitting a fixed offset away from it, which makes measured minus ceiling the error of that filter
rather than a statement about occupancy. Each CUSUM arm was therefore driven by a near-zero-mean
statistic, exactly the i.i.d. regime Siegmund's ARL$_0$ assumes, and over 120 room-weeks it duly
fired. The trace is unambiguous: in a ten-seat room whose design-occupancy trajectory settles near
1500 ppm, the ceiling sat 50 to 80 ppm \emph{above} the reading all day, so nine people arriving
walked the measurement straight through it. The floor had the mirror defect, dragged \emph{up} onto
the reading and crossed from above by every departure, which is why the lower arm's confirmations
cluster at 16:00 to 19:00 when offices empty. And the cascade is the same defect at one remove: a
false \texttt{STUCK\_DAMPER} pins the room's believed opening in \texttt{AhuFlowModel}, the split is
a ratio over all six rooms, so the other five are credited with air they are not getting, their
filtered ceilings fall below their readings and they confirm in turn. One seed latched four rooms
in 53 minutes and another latched all six (ADR-0010). The cross-room cascade an earlier pass
blamed, which charged the \emph{response} to a diagnosis rather than the belief behind it, stays
ruled out \emph{on the drawn weekday}: T-34b records 0 fault events and 0 degraded rooms across all
270 drawn-weekday observations in rooms with no injected fault. \textbf{At design occupancy it is
not ruled out at all}: the same experiment's second arm accuses 215 of 270 observed rooms, and
the supervisor-dead pass shows the accusations are the monitor's rather than the compensation's.
That is this same rate, measured by a different experiment in a regime R-3's own replay does not
visit, and §\ref{sec:11.6} records it as the cross-requirement pattern it is.

\textbf{Half of that defect is corrected and half is deliberately not, and the half that is not is
the more interesting decision.} Taken: the guarded lower arm now runs against the one bound that
needs no anchor (the outdoor concentration less the overnight-floor allowance), and a confirmation
that pins a damper re-anchors the other five rooms and resets their arms, because a step change in
the believed airflow makes every ppm banked before it a statement about a different building. That
is what took the rate from 1.30 to 0.55 and the clean weeks from 9 to 16, and it moved \emph{no}
other requirement: R-1's stepped p95 of 6 min and closed-loop 15 of 20, R-2, R-4 at 23 of 23, R-6
and R-7's drawn-weekday 0 of 270 are all unchanged. \textbf{Not taken: making the upper arm's ceiling a bound as
well.} That variant was built and measured, and it makes \textbf{R-3 pass outright}: 0.00 false
confirmations per week, 20 of 20 weeks clean. It was \textbf{declined on evidence}, because it
takes \textbf{R-1's closed-loop detection from 15 of 20 to 1 of 20}, \textbf{T-27 from 23 of 23
assertions to 19} and \textbf{R-7 from 0 of 270 observed rooms failing to 2}, and because sweeping
the plausibility ceiling's headroom finds no setting that recovers them: detection is 1 of 20 at
1.00 and 1 of 20 at 1.25, and below 1.00 the ceiling stops bounding occupancy at all. The
arithmetic behind that is structural rather than a threshold. The closed-loop fault is worth about
\textbf{86 ppm} of steady state at A124's drawn occupancy of 4 to 5 of its 8 desks, against
\textbf{270 to 600 ppm} of ceiling headroom; and those detections are real rather than the same
false alarms landing on a faulted room, because run as a pair the baseline latches A124 15 of 20
with the fault injected and \textbf{0 of 20 with nothing injected}. Both numbers come from one
knob, what the CUSUM's reference is, and there is no setting of it that meets R-3 and keeps R-1's
closed-loop arm. Breaking that tie needs a second, occupancy-free detection path rather than a
better reference (\texttt{ZoneObserver.stepAmplitudeRatio}, which the classifier already reads as
a signal and no CUSUM does), and that is design work rather than a parameter (ADR-0010). What is
left unattributed is narrower than it was and worth stating exactly: the residual \textbf{11}
confirmations sit in 4 weeks (6 on seed 2, 3 on seed 20, and one each on seeds 7 and 11), and seed
2's six are all on the upper arm, all after one primary on A124 at Wednesday 13:39.

\textbf{R-6 is met, and its headline number is an extrapolation rather than a measurement.} The
sweep runs to \textbf{25\,600 zones at a p95 of about 6 s}, roughly a tenth of the 60 s interval and
a factor of over 4000 above the six rooms the system actually serves, so no measured point breaks the
budget. That largest point is quoted to the precision the machine supports rather than to three
decimals: it has come out at 5.676, 5.740, 5.759, \textbf{6.008} and 6.052 s over five runs on the
same laptop (the fourth of those being the committed run this report quotes from), a spread of
about 7 \%, because this is the one experiment in the harness that measures real time. The
breaking point is therefore read off a fit: over the nine points from 100 zones up,
p95 = $1.196\times10^{-8}\,n^{1.97}$ s with $r^2$ = 0.9991, which crosses 60 s at about
\textbf{86\,000 zones}. And it moves with the run that fitted it: the committed run gives 85\,992,
while earlier runs of the same code gave 87\,338, 88\,109 and 87\,521. So the figure carries one
significant digit, 86\,000 is the committed one, and the 86\,000–88\,000 spread is its only error
bar. The harness labels it an extrapolation in those words and this
report keeps the label; what is \emph{measured} is the shape and the cause. The fitted exponent is 1.97 against a
theoretical 2 and reproduces across runs, and the direct evidence is the $p95/n^2$ column: it falls
while per-zone fixed costs dominate and then settles on 8.2–9.2 ns across the last seven points, 400
to 25\,600 zones, a 64-fold range of $n$, which neither a linear nor an $n \log n$ cost would do. The
saturated element is named rather than guessed: \texttt{AhuFlowModel.believedFlowTo} sums every
zone's opening to normalise the proportional split, and every zone's observer asks it, so one cycle
is $\Theta(n^2)$ in the zone count. That is the computational price of ADR-0008 option B: one
shared handler means every room's prediction depends on every other room's damper, and there is no
decomposition that avoids the sum. At six rooms the price is 0.007 ms, which is why the decision was
right. Reading ingestion was never close to the bottleneck: it stays at 0.89 ms p95 for 25\,600
readings, three orders of magnitude under the cycle, against a risk register that had nominated the
time-series query as the likely limit.

\textbf{R-6 now has a second ceiling, and it is three orders of magnitude below the first.} The
capability is no longer harness-only: \texttt{ZONE\_COUNT} is settable from the environment, all
\factAuthoredServices{} project-authored containers read the same one the way they already share \texttt{REAL\_EPOCH}
and \texttt{TIME\_SCALE}, and the stack was brought up and observed at 12 and at 58 zones with the
default six-room corridor unchanged end to end. \textbf{58 is a hard ceiling, and it is BuildSim's
rather than the monitor's}: it refuses to host equipment in a room its floorplan does not have, so a
deployable zone must be named after a real one, and level 0 of the A building has 52 further rooms
of comparable size. Verified in both directions: \texttt{ZONE\_COUNT=58} runs, \texttt{ZONE\_COUNT=60}
dies with \texttt{room "A124S9" does not exist on level "level0"}. So \textbf{R-6's breaking point
remains a harness measurement} and this report does not claim otherwise; what the wiring buys is a
deployed system that can be scaled and watched at all. One honesty note travels with the synthetic
zones: a zone past the sixth takes a real level-0 room's \emph{name and location} so BuildSim can
carry its sensor and damper, while its volume, floor area, damper authority and occupancy schedule
are a cycled copy of one of the six. \texttt{A119} is 109.6~m² in the floorplan and is simulated as
a copy of \texttt{A125A}'s 18~m². A synthetic zone is therefore representative in size, time constant
and head count, and is not a claim about the room whose name it borrows.

\textbf{R-7 is met on the drawn weekday and fails at design occupancy, and the two numbers must
always travel together.} Each of the six rooms takes its turn as the injected room, on three seeds,
in three configurations (the deployed guarded lower arm with the supervisor live, the arm off with
the supervisor live, and the guarded arm with the supervisor dead, 90 observed rooms each), and each
observed room is compared with \emph{itself} in a fault-free run on the same seed rather than with a
neighbour. On the \textbf{drawn weekday}, across all three passes, \textbf{0 of 270} observed rooms
failed: no fault event in a room with no injected fault, no room degraded, no room outside the
fault-free CO₂ band, and both the plan's own airflow figure and an independently recomputed
\texttt{FlowSplit} at 0.0000 against a 5 \% tolerance.

\textbf{At design occupancy the same experiment fails 231 of its 270 observed rooms}, and the way it
fails is the result rather than the count. The two failing clauses are \textbf{disjoint}: 215 rooms
were themselves \emph{accused of a fault they did not have}, and a further 16 sat outside the
airflow band at a worst deviation of \textbf{0.058919}: 215 plus 16 is the 231, so neither clause
is standing in for the other. They also attribute in opposite directions, and the supervisor-dead
pass is what separates them. Accusations run \textbf{85 of 90} rooms with the supervisor dead against
130 of 180 with it live, so they are the regime and the monitor rather than the response to the
fault: compensation slightly \emph{suppresses} them. Airflow breaches run \textbf{0 of 90} dead
against 16 of 180 live, so that half is the compensation's own cost and nothing else's: a dead
supervisor produces no plan to charge anyone for. \textbf{The accusing half is R-3's false-alarm
rate arriving inside R-7's experiment}, which is the same cross-requirement pattern §\ref{sec:11.2} already
records for R-2, and ADR-0009 predicted it in those words while noting that no experiment asserted
R-7 in that regime. One does now, and they collide.

\textbf{The mode clause is no longer unexercised, and the requirement now fails under either
reading.} The drawn-weekday pass reported it as satisfied by accident rather than by design: a plan
was in force in 16 of the 18 injections in each supervised pass and every one came back with a trim
\emph{factor} of exactly 1.00 — no trim at all, which publishes no damper command, so no healthy
room left \texttt{NORMAL} control, and the
record named the honest re-test as an experiment that \emph{forces} a trim. The design-occupancy arm
is it: \textbf{40 of 540} rooms reach \texttt{COMPENSATING}. So R-7's literal wording fails on 40
rooms and ADR-0008 consequence 9's defensible reading fails on 16, and choosing between the two
readings no longer rescues the requirement. \texttt{DEGRADED} reaches \textbf{130 of 540}, which is
the half that matters: a degraded room is one the monitor accused and took off demand control, and
§\ref{sec:11.2} records 130 of the 215 falsely latched rooms in exactly that state.

One of the two airflow columns cannot falsify this requirement and should stop being quoted beside
one that can. \texttt{CompensationPlan.worstHealthyFlowDeviation()} reaches 0.0500 and no further in
any regime, because the solver caps the fan at the speed that makes the figure exactly
\texttt{HEALTHY\_FLOW\_TOLERANCE} and then reports that cap: it is the solver quoting its own
constraint back. The independently recomputed \texttt{FlowSplit} deviation is the only one of the two
that can fail, and it does. Behind that sits an unresolved question that is ADR-level rather than a
code change: \textbf{the 5 \% bound's reference plant is ambiguous}. The solver bounds its claim over
each latched room's \emph{latch-time} opening while the recomputation uses its \emph{pre-latch} one,
which on one seed gives the same room 48.7940 against 38.4015 and the ratios 1.050000 against
1.058919. The two coincide until a second room latches inside one episode, which the drawn weekday
never does and design occupancy routinely does. Both are defensible readings of ``before the
compensation''; ADR-0008 consequence 9 never says which.

The one large column is the \emph{measured} end-to-end flow gap, \textbf{0.00 to 6.10 times the
fault-free flow} over the 540 observed-room rows, and it is not a violation: a damper frozen near its
minimum stops competing for the shared fan, so the five rooms \emph{gain} the air it is not taking.
ADR-0008 priced that coupling explicitly, and the 5 \% tolerance bounds what compensation takes from
the healthy rooms, not what the fault gives them.
